\section{Local smoothing laws at regular and singular minimizers}
\label{sec:rates}

Write
\[
 A(\lambda)=G+\lambda\Theta,
 \qquad s=\lambda-\lambda^\star.
\]
All statements in this section are local asymptotics at a fixed instance
under the printed nondegeneracy assumptions.  The strict rank-deficient law
below permits arbitrary shape and corank.  Only the boundary and crossover
laws use a transverse square corank-one crossing, because their proof isolates
one signed Rellich-analytic small mode.

\begin{lemma}[Exact analytic separation of the small singular cluster]
\label{lem:analytic-cluster-block-main}
Let \(A(s)=A^\star+s\Theta\), where \(\operatorname{rank}A^\star=r\).
Locally there are real-analytic orthogonal frames
\(U(s)=[U_+(s)\ U_0(s)]\) and \(V(s)=[V_+(s)\ V_0(s)]\), agreeing
at zero with fixed compact and null frames, such that
\[
 U(s)^\top A(s)V(s)=
 \begin{pmatrix}A_+(s)&0\\0&B(s)\end{pmatrix},
 \qquad B(s)=sC+s^2E(s),\quad C=U_0^\top\Theta V_0,
\]
where \(A_+(s)\in\mathbb R^{r\times r}\) is invertible and \(E\) is
real analytic.  Empty blocks are allowed.  If \(q=\min\{m,n\}\) and
singular values are listed in decreasing order, then
\[
 \sigma_{r+i}(A(s))=|s|\sigma_i(C)+O(s^2),
 \qquad 1\le i\le q-r,
\]
including the zero singular values of \(C\).
\end{lemma}

\begin{proof}
When \(r=0\), choose constant frames; the conclusion is immediate.
Otherwise a fixed spectral gap separates the zero and positive eigenvalue
clusters of \(A^\star{}^\top A^\star\) and \(A^\star A^\star{}^\top\).
The corresponding small-cluster Riesz projections \(\Pi_V(s),\Pi_U(s)\)
are analytic.  For a common contour enclosing zero and no positive
eigenvalue, the resolvent identity
\[
 A(s)(zI-A(s)^\top A(s))^{-1}
 =(zI-A(s)A(s)^\top)^{-1}A(s)
\]
implies \(A(s)\Pi_V(s)=\Pi_U(s)A(s)\), including unequal left and
right nullities.  Analytic orthonormal frames for these projections and
their complements give the exact block decomposition.  Such frames can
be constructed by projecting the fixed frames and multiplying by the
inverse square root of their Gram matrices.  Since
\(A^\star V_0=0\) and \(U_0^\top A^\star=0\), differentiation of
\(B(s)=U_0(s)^\top A(s)V_0(s)\) gives \(B(0)=0\) and \(B'(0)=C\);
analytic division gives the stated \(E\).  Finally, the small singular
values of \(A(s)\) are exactly those of \(B(s)\).  The singular-value
perturbation bound applied to \(B(s)-sC=O(s^2)\) proves the indexed formula.
No choice of individual singular vectors within a repeated cluster is used.
\end{proof}

\subsection{Full-rank regular minimizers}

\begin{theorem}[Quadratic regular displacement]
\label{thm:rate-regular}
Assume that \(\lambda^\star\) minimizes \(\phi\), that
\(A^\star=A(\lambda^\star)\) has full rectangular rank, and that
\(\kappa=\phi''(\lambda^\star)>0\).  If
\(A^\star=U\Sigma V^\top\), then
\begin{equation}
 \lambda_\delta-\lambda^\star
 =-\frac{b^\star}{\kappa}\delta^2+o(\delta^2),
 \qquad
 b^\star=-\frac12
 \langle\Theta,U\Sigma^{-2}V^\top\rangle.
 \label{eq:regular-rate-main}
\end{equation}
The matrix in the pairing is basis independent:
\begin{equation}
 U\Sigma^{-2}V^\top=
 \begin{cases}
 A^\star((A^\star)^\top A^\star)^{-3/2},&m\ge n,\\
 (A^\star(A^\star)^\top)^{-3/2}A^\star,&m<n.
 \end{cases}
 \label{eq:regular-intrinsic-main}
\end{equation}
Thus the displacement is \(O(\delta^2)\), and it is
\(\Theta(\delta^2)\) when \(b^\star\ne0\).
\end{theorem}

\begin{proof}
In a full-rank neighborhood,
\[
 \frac{\sigma}{\sqrt{\sigma^2+\delta^2}}
 =1-\frac{\delta^2}{2\sigma^2}+O(\delta^4)
\]
uniformly.  Functional calculus therefore gives
\[
 Z_\delta(A(\lambda))
 =\Sc(A(\lambda))
 -\frac{\delta^2}{2}
 U(\lambda)\Sigma(\lambda)^{-2}V(\lambda)^\top
 +O(\delta^4).
\]
The displayed matrix functions are analytic and basis independent even at
repeated positive singular values.  Pairing with \(\Theta\) yields
\(h_\delta(\lambda)=\phi'(\lambda)+\delta^2b(\lambda)+O(\delta^4)\), with
\(b(\lambda^\star)=b^\star\).  The scalar implicit function theorem at
\(\phi'(\lambda^\star)=0\), \(\phi''(\lambda^\star)=\kappa\) proves
\eqref{eq:regular-rate-main}.
\end{proof}

\subsection{Strict rank-deficient displacement at arbitrary corank}

\begin{theorem}[Strict displacement at arbitrary corank]
\label{thm:rate-strict-general}
Let \(\lambda^\star\) minimize \(\phi\), let
\(A^\star=A(\lambda^\star)\) be rank deficient, and define \(a,C,\beta\)
by \eqref{eq:a-C-beta-main}.  Assume the strict condition
\(|a|<\beta\).  Let
\(\sigma_1\ge\cdots\ge\sigma_{r_C}>0\) be the positive singular values
of \(C\), and let \(\eta_R\ge0\) be the unique solution of
\begin{equation}
 \sum_{i=1}^{r_C}
 \frac{\eta_R\sigma_i^2}{\sqrt{1+\eta_R^2\sigma_i^2}}=|a|.
 \label{eq:eta-strict-rate-main}
\end{equation}
Then
\begin{equation}
 \frac{\lambda_\delta-\lambda^\star}{\delta}
 \longrightarrow-\sign(a)\eta_R.
 \label{eq:strict-general-rate-main}
\end{equation}
Moreover, with \(q=\min\{m,n\}\),
\begin{equation}
 |\lambda_\delta-\lambda^\star|
 \le \frac{q}{\beta-|a|}\,\delta.
 \label{eq:strict-general-bound-main}
\end{equation}
No square-shape, corank-one, or rank-one-normal assumption is required.
The displacement is \(\Theta(\delta)\) exactly when \(a\ne0\); when
\(a=0\), it is \(o(\delta)\).
\end{theorem}

\begin{proof}
Put \(q=\min\{m,n\}\), \(r=\rank A^\star\), and
\(t_\delta=(\lambda_\delta-\lambda^\star)/\delta\).  Strict optimality
and \cref{thm:interval} give the global sharp-growth estimate
\begin{equation}
 \phi(\lambda)\ge \phi(\lambda^\star)
 +(\beta-|a|)|\lambda-\lambda^\star|.
 \label{eq:strict-sharp-growth-main}
\end{equation}
Combining it with \eqref{eq:smoothing-uniform} and optimality of
\(\lambda_\delta\) yields \eqref{eq:strict-general-bound-main}.
Thus \((t_\delta)\) is bounded before any local expansion is used.

Define the rescaled convex functions
\[
 F_\delta(t)=
 \frac{\Psi_\delta(A(\lambda^\star+\delta t))
       -\Psi_\delta(A^\star)}{\delta}.
\]
Choose a spectral gap separating the \(r\) positive singular values of
\(A^\star\) from zero.  The sum \(\ell(s)\) of their analytic separated
cluster is smooth even when singular values inside that cluster cross, and
\(\ell'(0)=\langle\Theta,U_rV_r^\top\rangle=a\).  Since
\(\psi_\delta(\sigma)=\sigma-\delta+O(\delta^2)\) uniformly on the
separated cluster, its contribution to \(F_\delta(t)\) converges to \(at\),
uniformly for \(t\) in compact sets.

For the small cluster, \cref{lem:analytic-cluster-block-main} gives,
uniformly near zero,
\[
 \sigma_{r+i}(A(\lambda^\star+s))
 =|s|\sigma_i(C)+O(s^2),
 \qquad 1\le i\le q-r,
\]
where zero singular values of \(C\) are included.  Setting \(s=\delta t\)
and using that \(x\mapsto\sqrt{1+x^2}-1\) is Lipschitz gives compact-uniform
convergence
\begin{equation}
 F_\delta(t)\longrightarrow
 F(t):=at+\sum_{i=1}^{r_C}
 \left(\sqrt{1+\sigma_i^2t^2}-1\right).
 \label{eq:strict-epi-limit-main}
\end{equation}

The function \(F\) is strictly convex, and
\(F(t)=at+\beta|t|+o(|t|)\); hence \(|a|<\beta\) makes it coercive.
Its unique minimizer solves
\[
 a+\sum_{i=1}^{r_C}
 \frac{\sigma_i^2t}{\sqrt{1+\sigma_i^2t^2}}=0,
\]
so it equals \(-\sign(a)\eta_R\).  The localization
\eqref{eq:strict-general-bound-main} puts every minimizer of \(F_\delta\) in
one fixed compact interval.  Compact-uniform convergence then implies that
every accumulation point minimizes \(F\); uniqueness proves
\eqref{eq:strict-general-rate-main}.  The final order statements follow
because \(\eta_R>0\) if and only if \(a\ne0\).
\end{proof}

\begin{remark}[One scalar governs both path and limit]
\label{rem:one-scalar-path-center}
Throughout this section the subscript \(R\) marks objects attached to the
pseudo-Huber regularizer, rather than a common tensor type.
The multiplier \(\eta_R\) in \eqref{eq:strict-general-rate-main} is exactly
the multiplier in the smooth certificate formula
\eqref{eq:eta-R-main}--\eqref{eq:KR-main}.  Thus the first displacement of
the smoothing path and the kernel block of its limiting certificate are two
manifestations of the same scalar balance.  More precisely,
\(K_R=Z_1(t_0C)\), where
\(t_0=-\operatorname{sign}(a)\eta_R\); the next theorem shows that this
identity is the zero-scale value of an analytic blow-up path.
\end{remark}

\begin{theorem}[First-order direction law at a strict minimizer]
\label{thm:direction-strict-general}
Under the hypotheses of \cref{thm:rate-strict-general}, put
\(t_0=-\operatorname{sign}(a)\eta_R\).  For the fixed instance there is a
jointly real-analytic blow-up extension \(\mathcal Z(\varepsilon,t)\), defined
near \((0,t_0)\), such that, for \(\varepsilon>0\),
\begin{equation}
 \mathcal Z(\varepsilon,t)
 =Z_\varepsilon(A(\lambda^\star+\varepsilon t)),
 \qquad
 \mathcal Z(0,t)=U_rV_r^\top+U_0Z_1(tC)V_0^\top.
 \label{eq:strict-blowup-map-main}
\end{equation}
Let
\[
 H(\varepsilon,t)=\langle\Theta,\mathcal Z(\varepsilon,t)\rangle,
 \qquad
 d_R:=\partial_tH(0,t_0)
 =\sum_{i=1}^{r_C}
 \frac{\sigma_i^2}{(1+t_0^2\sigma_i^2)^{3/2}}>0,
\]
and set
\begin{equation}
 t_1=-\frac{\partial_\varepsilon H(0,t_0)}{d_R},
 \qquad
 \mathcal M_R=\partial_\varepsilon\mathcal Z(0,t_0)
 -\frac{\langle\Theta,\partial_\varepsilon\mathcal Z(0,t_0)\rangle}{d_R}
  \partial_t\mathcal Z(0,t_0).
 \label{eq:strict-direction-constant-main}
\end{equation}
Then
\begin{align}
 \lambda_\delta-\lambda^\star
 &=\delta t_0+\delta^2t_1+O(\delta^3),
 \label{eq:strict-multiplier-second-main}\\
 \Phi_\delta
 &=\Phi_R^\star+\delta\mathcal M_R+O(\delta^2).
 \label{eq:strict-direction-first-main}
\end{align}
Consequently the returned-direction error is \(O(\delta)\), and it is
asymptotic to \(\delta\|\mathcal M_R\|_F\) when
\(\mathcal M_R\ne0\).  If \(\mathcal M_R=0\), it is \(O(\delta^2)\).
In particular, \(a=0\) forces \(t_0=t_1=0\) and
\(\mathcal M_R=0\), so that
\begin{equation}
 \lambda_\delta-\lambda^\star=O(\delta^3),
 \qquad
 \|\Phi_\delta-\Phi_R^\star\|_F=O(\delta^2).
 \label{eq:strict-zero-superconvergence-main}
\end{equation}
The constants are local to the fixed instance; no uniformity as the spectral
gap or \(\beta-|a|\) tends to zero is asserted.
\end{theorem}

\begin{proof}
Let \(r=\operatorname{rank}A^\star\).  Apply
\cref{lem:analytic-cluster-block-main} to obtain local analytic orthogonal frames
\(U(s)=[U_+(s)\ U_0(s)]\) and \(V(s)=[V_+(s)\ V_0(s)]\) for which
\[
 U(s)^\top A(\lambda^\star+s)V(s)
 =\begin{pmatrix}A_+(s)&0\\0&B(s)\end{pmatrix}.
\]
Here \(A_+(s)\in\mathbb R^{r\times r}\) is invertible,
\(B(s)\in\mathbb R^{(m-r)\times(n-r)}\), and
\(B(0)=0\), \(B'(0)=C\).  Analytic division gives
\(B(s)=s\widetilde B(s)\), with \(\widetilde B(0)=C\); unequal left and
right nullities and zero singular values of \(C\) cause no difficulty.
The frames are chosen to agree at \(s=0\) with the fixed compact and null
frames used to define \(a\) and \(C\).

Write
\[
 \mathcal E_s(K)=U_0(s)KV_0(s)^\top,
 \qquad
 P_\varepsilon(s)=
 U_+(s)Z_\varepsilon(A_+(s))V_+(s)^\top.
\]
Orthogonal equivariance, block separation, and the scaling identity
\(Z_\varepsilon(B)=Z_1(B/\varepsilon)\) give, for \(\varepsilon>0\),
\begin{equation}
 Z_\varepsilon(A(\lambda^\star+\varepsilon t))
 =P_\varepsilon(\varepsilon t)
 +\mathcal E_{\varepsilon t}
  \left[Z_1\!\left(t\widetilde B(\varepsilon t)\right)\right].
 \label{eq:strict-blowup-block-main}
\end{equation}
The positive block is analytic in \((s,\varepsilon^2)\), while \(Z_1\) is
real analytic on the entire rectangular matrix space.  Thus the right-hand
side extends jointly analytically across \(\varepsilon=0\) and proves
\eqref{eq:strict-blowup-map-main}.  This construction uses analytic cluster
subspaces, not analytic individual singular vectors, so repeated positive
singular values are allowed \citep[Chapter II]{kato1995perturbation}.

At zero scale,
\[
 H(0,t)=a+\langle C,Z_1(tC)\rangle
 =a+\sum_{i=1}^{r_C}
   \frac{t\sigma_i^2}{\sqrt{1+t^2\sigma_i^2}}.
\]
Its unique zero is \(t_0\), and differentiation gives the stated positive
constant \(d_R\).  The analytic implicit function theorem therefore gives
\(t(\varepsilon)=t_0+\varepsilon t_1+O(\varepsilon^2)\).  For positive
\(\varepsilon\), strict convexity and
\cref{thm:rate-strict-general} identify this local branch with
\((\lambda_\varepsilon-\lambda^\star)/\varepsilon\).
Since \(Z_1(t_0C)=K_R\), Taylor expansion of
\(\mathcal Z(\delta,t(\delta))\) proves
\eqref{eq:strict-multiplier-second-main}--\eqref{eq:strict-direction-first-main}
and \eqref{eq:strict-direction-constant-main}.  When \(a=0\), one has
\(t_0=0\), while
\(Z_\varepsilon(A^\star)=\Sc(A^\star)+O(\varepsilon^2)\).  Hence
\(\partial_\varepsilon H(0,0)=0\) and
\(\partial_\varepsilon\mathcal Z(0,0)=0\), proving
\eqref{eq:strict-zero-superconvergence-main}.
\end{proof}

\begin{proposition}[Intrinsic computable first-correction formula]
\label{prop:strict-constants-closed-main}
Use the analytic cluster frames in the proof of
\cref{thm:direction-strict-general}, write
\(\widetilde B(s)=B(s)/s\) with its analytic value
\(\widetilde B(0)=C\), and define the basis-free frozen-scale field
\begin{equation}
 \mathcal W(s,t)=P_{\rm big}(s)
 +U_0(s)Z_1\!\bigl(t\widetilde B(s)\bigr)V_0(s)^\top,
 \qquad
 P_{\rm big}(s)=U_+(s)\Sc(A_+(s))V_+(s)^\top.
 \label{eq:frozen-field-main}
\end{equation}
Then \(\mathcal W(0,t_0)=\Phi_R^\star\) and
\begin{equation}
 \partial_\varepsilon\mathcal Z(0,t_0)
 =t_0\partial_s\mathcal W(0,t_0).
 \label{eq:frozen-chain-main}
\end{equation}
Consequently the constants in
\eqref{eq:strict-direction-constant-main} are
\begin{align}
 t_1&=-\frac{t_0
 \langle\Theta,\partial_s\mathcal W(0,t_0)\rangle}{d_R},
 \label{eq:t1-closed-main}\\
 \mathcal M_R&=t_0\left[
 \partial_s\mathcal W(0,t_0)
 -\frac{\langle\Theta,\partial_s\mathcal W(0,t_0)\rangle}{d_R}
  \partial_t\mathcal W(0,t_0)
 \right].
 \label{eq:MR-closed-main}
\end{align}
In particular, \(t_0=0\) makes both constants vanish.  If \(t_0\ne0\),
then
\begin{align*}
 \mathcal M_R=0
 &\iff \partial_s\mathcal W(0,t_0)
       \in\operatorname{span}\{\partial_t\mathcal W(0,t_0)\},\\
 t_1=0
 &\iff \langle\Theta,\partial_s\mathcal W(0,t_0)\rangle=0,
\end{align*}
and both vanish if and only if
\(\partial_s\mathcal W(0,t_0)=0\).
\end{proposition}

\begin{proof}
Both terms in \eqref{eq:frozen-field-main} are invariant under analytic
orthogonal changes of the left and right cluster frames, by orthogonal
equivariance of \(Z_1\).  In \eqref{eq:strict-blowup-block-main}, the positive
block satisfies
\(P_\varepsilon(s)=P_{\rm big}(s)+O(\varepsilon^2)\), jointly analytically
in \((s,\varepsilon^2)\).  All remaining \(\varepsilon\)-dependence enters
through \(s=\varepsilon t\).  Hence
\[
 \mathcal Z(\varepsilon,t)
 =\mathcal W(\varepsilon t,t)+O(\varepsilon^2),
\]
and the chain rule proves \eqref{eq:frozen-chain-main}.  Substitution into
\eqref{eq:strict-direction-constant-main} gives
\eqref{eq:t1-closed-main}--\eqref{eq:MR-closed-main}.  Finally,
\(\langle\Theta,\partial_t\mathcal W(0,t_0)\rangle=d_R>0\);
the stated vanishing alternatives follow directly.  Notice in particular
that parallelism alone cancels \(\mathcal M_R\), but need not cancel \(t_1\).
The full-rank polar-block derivatives in this formula may be evaluated by
standard Fr\'echet-derivative algorithms \citep{gawlik2017frechet}; the new
feature here is their coupling to the analytic rank-deficient blow-up.
\end{proof}

\begin{corollary}[Sharp compact-root superconvergence]
\label{cor:compact-root-superconvergence-main}
Under \cref{thm:direction-strict-general}, suppose \(a=0\).  By
\cref{thm:interval}, this is exactly compact-polar stationarity at
\(\lambda^\star\).  Let
\begin{equation}
 J_\star=U_r\Sigma_r^{-2}V_r^\top,
 \qquad d_0=\|C\|_F^2,
 \qquad \chi_\star=\langle\Theta,J_\star\rangle.
 \label{eq:compact-root-coefficients-main}
\end{equation}
where \(J_\star=0\) when the positive singular block is empty.  Strictness
gives \(C\ne0\), hence \(d_0>0\).
Then \(\Phi_R^\star=\Sc(A^\star)\), and
\begin{align}
 \lambda_\delta-\lambda^\star
 &=\frac{\chi_\star}{2d_0}\,\delta^3+O(\delta^4),
 \label{eq:compact-root-cubic-main}\\
 \Phi_\delta-\Phi_R^\star
 &=\delta^2\mathcal N_0+O(\delta^3),
 \qquad
 \mathcal N_0=-\frac12J_\star
 +\frac{\chi_\star}{2d_0}U_0CV_0^\top.
 \label{eq:compact-root-quadratic-main}
\end{align}
Moreover,
\begin{equation}
 \|\mathcal N_0\|_F^2
 =\frac14\|J_\star\|_F^2+\frac{\chi_\star^2}{4d_0}.
 \label{eq:compact-root-direction-norm-main}
\end{equation}
Thus the multiplier error is exactly cubic when \(\chi_\star\ne0\), and the
direction error is exactly quadratic whenever \(\rank A^\star>0\).  If
\(\chi_\star=0\), put
\begin{equation}
 J_2=U_r\Sigma_r^{-4}V_r^\top,
 \qquad \chi_2=\langle\Theta,J_2\rangle .
 \label{eq:compact-root-second-coefficients-main}
\end{equation}
Then the next coefficients are
\begin{align}
 \lambda_\delta-\lambda^\star
 &=-\frac{3\chi_2}{8d_0}\,\delta^5+O(\delta^6),
 \label{eq:compact-root-quintic-main}\\
 \Phi_\delta-\Phi_R^\star
 &=-\frac{\delta^2}{2}J_\star
 +\delta^4\left[\frac38J_2
 -\frac{3\chi_2}{8d_0}U_0CV_0^\top\right]+O(\delta^5).
 \label{eq:compact-root-quartic-main}
\end{align}
In particular, the multiplier error is exactly quintic if \(\chi_2\ne0\).
\end{corollary}

\begin{proof}
Here \(t_0=0\), \(d_R=d_0>0\), and
\[
 \mathcal Z(\varepsilon,0)=Z_\varepsilon(A^\star)
 =\Sc(A^\star)-\frac{\varepsilon^2}{2}J_\star+O(\varepsilon^4),
 \qquad
 \partial_t\mathcal Z(0,0)=U_0CV_0^\top.
\]
It follows that
\(H(\varepsilon,0)=-\chi_\star\varepsilon^2/2+O(\varepsilon^4)\).
The analytic implicit solution therefore satisfies
\[
 t(\varepsilon)=\frac{\chi_\star}{2d_0}\varepsilon^2
 +O(\varepsilon^3).
\]
Since \(\lambda_\varepsilon-\lambda^\star=\varepsilon t(\varepsilon)\),
substitution into the analytic blow-up gives
\eqref{eq:compact-root-cubic-main}--\eqref{eq:compact-root-quadratic-main}.
The range--range and null--null blocks in \(\mathcal N_0\) are Frobenius
orthogonal, which proves \eqref{eq:compact-root-direction-norm-main}.
If \(\chi_\star=0\), the separated positive block has the uniform expansion
\[
 \frac{\sigma}{\sqrt{\sigma^2+\varepsilon^2}}
 =1-\frac{\varepsilon^2}{2\sigma^2}
  +\frac{3\varepsilon^4}{8\sigma^4}+O(\varepsilon^6).
\]
Consequently
\(H(\varepsilon,0)=3\chi_2\varepsilon^4/8+O(\varepsilon^6)\), and the
implicit branch satisfies
\(t(\varepsilon)=-3\chi_2\varepsilon^4/(8d_0)+O(\varepsilon^5)\).
This proves \eqref{eq:compact-root-quintic-main}.  Now
\(s=\varepsilon t(\varepsilon)=O(\varepsilon^5)\).  The quartic term of
the positive block and
\(t(\varepsilon)\partial_t\mathcal Z(0,0)\), with
\(\partial_t\mathcal Z(0,0)=U_0CV_0^\top\), give
\eqref{eq:compact-root-quartic-main}.  Indeed, analyticity yields
\(\mathcal Z(\varepsilon,t(\varepsilon))
=\mathcal Z(\varepsilon,0)+t(\varepsilon)\partial_t\mathcal Z(0,0)
+O(\varepsilon|t(\varepsilon)|+t(\varepsilon)^2)\), so every remaining
term is \(O(\varepsilon^5)\).
\end{proof}

\begin{remark}[An odd-order multiplier hierarchy]
\label{rem:compact-root-hierarchy-main}
The cubic and quintic laws are the first two levels of a complete multiplier
hierarchy.  For an integer \(k\ge1\), let
\[
 b_k=(-1)^k4^{-k}\binom{2k}{k},\qquad
 J_k=U_r\Sigma_r^{-2k}V_r^\top,\qquad
 \chi_k=\langle\Theta,J_k\rangle.
\]
If an integer \(p\ge1\) satisfies
\(a=0\) and \(\chi_1=\cdots=\chi_{p-1}=0\), then the binomial expansion
used above and the analytic implicit function theorem give
\begin{equation}
 \lambda_\delta-\lambda^\star
 =-\frac{b_p\chi_p}{d_0}\,\delta^{2p+1}
  +O(\delta^{2p+2}).
 \label{eq:compact-root-hierarchy-main}
\end{equation}
Thus each further orthogonality cancellation raises the multiplier by two
orders; \(p=1\) is \eqref{eq:compact-root-cubic-main} and \(p=2\) is
\eqref{eq:compact-root-quintic-main}.  This hierarchy is local to the fixed
strict crossing and does not assert that the cancellations are generic.
\end{remark}

\begin{corollary}[Finite cancellation criterion for an exact multiplier]
\label{cor:finite-cancellation-main}
Under \cref{thm:direction-strict-general}, suppose \(a=0\), and let
\(d\) be the number of distinct positive singular values of \(A^\star\).
If \(d=0\), then \(\lambda_\delta=\lambda^\star\) for every \(\delta>0\).
If \(d\ge1\), the following statements are equivalent:
\begin{enumerate}
 \item \(\chi_1=\cdots=\chi_{d-1}=0\), with no conditions when \(d=1\);
 \item \(\langle\Theta,U_jV_j^\top\rangle=0\) for each positive
 singular-value level \(j\), where \(U_j,V_j\) contain its paired frames;
 \item \(\lambda_\delta=\lambda^\star\) for every \(\delta>0\).
\end{enumerate}
Otherwise, the first nonzero \(\chi_p\) occurs with \(1\le p\le d-1\),
and \eqref{eq:compact-root-hierarchy-main} gives its nonzero leading bias.
\end{corollary}

\begin{proof}
If \(d=0\), then \(A^\star=0\) and \(h_\delta(\lambda^\star)=0\).
For \(d\ge1\), write the distinct positive levels as \(\nu_j\), put
\(w_j=\langle\Theta,U_jV_j^\top\rangle\), and set \(x_j=\nu_j^{-2}\).
Then \(a=\sum_jw_j=0\) and \(\chi_k=\sum_jw_jx_j^k\).  The Vandermonde
matrix with rows \(k=0,\ldots,d-1\) is invertible, proving equivalence of
the first two statements.  Moreover,
\[
 h_\delta(\lambda^\star)
 =\sum_{j=1}^d\frac{w_j}{\sqrt{1+\delta^2x_j}}.
\]
Thus vanishing weights give an exact root for every \(\delta>0\), and
strict convexity identifies the unique multiplier.  Conversely, a root
for every positive \(\delta\) makes all Taylor coefficients at zero
vanish, in particular the first \(d\) moments.  This proves the equivalence
and the bound on the first nonzero moment.  The group polar matrices
\(U_jV_j^\top\) are independent of the paired frame chosen within a level.
\end{proof}

At a fixed strict rank-deficient minimizer, these results give a sharp
dichotomy.  Compact-polar stationarity holds exactly when \(a=0\), in which
case smoothing has \(O(\delta^3)\) multiplier bias and \(O(\delta^2)\)
direction bias with
the refinements above.  If \(a\ne0\), compact-polar stationarity fails and the
multiplier bias is necessarily \(\Theta(\delta)\); the direction bias is
\(\Theta(\delta)\) when \(\mathcal M_R\ne0\), but may cancel to second order.

\subsection{A transverse corank-one normal form for boundary and crossover}

For the remaining results assume \(m=n=q\),
\(\operatorname{rank}A^\star=q-1\), and let \(u_0,v_0\) be unit left and
right null vectors.  Put
\begin{equation}
 \alpha=u_0^\top\Theta v_0,
 \qquad \beta=|\alpha|,
 \qquad \alpha\ne0.
 \label{eq:alpha-beta-main}
\end{equation}

\begin{lemma}[Analytic kink normal form]
\label{lem:kink-normal-form}
There are real-analytic functions \(\tau,\ell\) near zero such that
\begin{equation}
 \tau(s)=\alpha s+O(s^2),
 \qquad
 \phi(\lambda^\star+s)=\ell(s)+|\tau(s)|.
 \label{eq:kink-normal-form}
\end{equation}
Moreover,
\begin{equation}
 h_\delta(\lambda^\star+s)
 =\ell_\delta'(s)
 +\frac{\tau(s)\tau'(s)}
       {\sqrt{\tau(s)^2+\delta^2}},
 \qquad
 \ell_\delta'(s)=\ell'(s)+O(\delta^2)
 \label{eq:kink-smooth-derivative}
\end{equation}
uniformly on a fixed small neighborhood, where \(\ell_\delta'(s)\) denotes the
contribution to \(h_\delta\) from the separated nonzero singular-value cluster.
With \(a=\ell'(0)\), the exact
subdifferential is \([a-\beta,a+\beta]\).
\end{lemma}

\begin{proof}
Consider the analytic symmetric dilation
\[
 H(s)=\begin{pmatrix}
 0&A^\star+s\Theta\\
 (A^\star+s\Theta)^\top&0
 \end{pmatrix}.
\]
Its zero eigenspace has basis \((u_0,0),(0,v_0)\), and the compression of
\(H'(0)\) is
\(\bigl(\begin{smallmatrix}0&\alpha\\\alpha&0\end{smallmatrix}\bigr)\).
Rellich perturbation theory \citep[Chapter II]{kato1995perturbation} therefore
gives a signed analytic pair \(\pm\tau(s)\) with
\(\tau'(0)=\alpha\).  The small singular value is \(|\tau(s)|\).  The sum of
the separated positive eigenvalue cluster of the dilation is analytic and
defines \(\ell(s)\), proving \eqref{eq:kink-normal-form}.  The nonzero cluster
stays uniformly separated from zero, so the regular expansion used in
\cref{thm:rate-regular} gives
\(\ell_\delta'=\ell'+O(\delta^2)\).  Differentiating the smoothed small mode
proves \eqref{eq:kink-smooth-derivative}; the two one-sided derivatives at
zero are \(a\mp\beta\).
\end{proof}

\subsection{The corank-one strict specialization and boundary kink}

\begin{corollary}[Corank-one strict displacement]
\label{thm:rate-strict}
Under \cref{lem:kink-normal-form}, suppose \(|a|<\beta\).  Then
\begin{equation}
 \frac{\lambda_\delta-\lambda^\star}{\delta}
 \longrightarrow
 t_0=-\frac{a}{\beta\sqrt{\beta^2-a^2}}.
 \label{eq:strict-rate-main}
\end{equation}
The displacement is \(O(\delta)\), and it is \(\Theta(\delta)\) when
\(a\ne0\).
\end{corollary}

\begin{proof}
Here \(C\) has the single positive singular value \(\beta\).
Equation \eqref{eq:eta-strict-rate-main} gives
\(\eta_R=|a|/[\beta\sqrt{\beta^2-a^2}]\), so
\cref{thm:rate-strict-general} reduces to \eqref{eq:strict-rate-main}.
\end{proof}

At the boundary, the small-mode expansion must be grouped with the complete
contact-side derivative; expanding the mode separately loses a linear term.
Here \(\phi'_{\rm con}\) denotes the derivative of the full-rank analytic branch
continued from the stated contact side.
The boundary law is not obtained by sending \(|a|\uparrow\beta\) in the
strict formula: \(\eta_R\to\infty\), so the \(\delta\)-scale balance itself
breaks down and is replaced by the \(\delta^{2/3}\) scale below.

\begin{theorem}[Boundary-kink two-thirds law]
\label{thm:rate-boundary}
Suppose \(|a|=\beta\), and choose the sign of the analytic branch in
\cref{lem:kink-normal-form} so that \(\alpha=\beta>0\).  Let the full-rank
contact side be \(s<0\) when \(a=\beta\) and \(s>0\) when \(a=-\beta\), and
let \(\phi_{\rm con}\) denote the analytic branch continued from that side.
Then \cref{lem:kink-normal-form} gives
\begin{equation}
 \phi'_{\rm con}(\lambda^\star+s)=\kappa s+O(s^2),
 \qquad
 \kappa:=
 \begin{cases}
  \ell''(0)-\tau''(0),&a=\beta,\\
  \ell''(0)+\tau''(0),&a=-\beta.
 \end{cases}
 \label{eq:contact-curvature-main}
\end{equation}
Convexity gives \(\kappa\ge0\); assume only the nondegeneracy
\(\kappa>0\).  Then
\begin{equation}
 \lambda_\delta-\lambda^\star
 =-\operatorname{sign}(a)
 \left(\frac{1}{2\beta\kappa}\right)^{1/3}
 \delta^{2/3}(1+o(1)).
 \label{eq:boundary-rate-main}
\end{equation}
\end{theorem}

\begin{proof}
The expansion in \eqref{eq:contact-curvature-main} is a consequence, not an
extra hypothesis.  On the contact side one has
\(\phi_{\rm con}=\ell-\tau\) when \(a=\beta\) and
\(\phi_{\rm con}=\ell+\tau\) when \(a=-\beta\).  Hence its derivative at
zero is respectively \(a-\beta=0\) or \(a+\beta=0\), and analyticity gives
the displayed expansion with the stated second derivative.

It suffices to treat \(a=\beta\); the other case follows after reversing
\(s\) and the signed analytic branch.  Since
\(\tau(0)=0\), \eqref{eq:kink-smooth-derivative} gives
\(h_\delta(\lambda^\star)=\ell_\delta'(0)=a+O(\delta^2)>0\) for all small
\(\delta\), and \(h_\delta\) is strictly increasing; hence the smooth root
lies on \(s<0\).
Put \(\chi(s)=\operatorname{sign}(\tau(s))\tau'(s)\).  Then
\(\chi(s)=-\beta+O(s)\) and
\(\tau(s)^2=\beta^2s^2(1+O(s))\) on that side.  In the regime
\(\delta/|s|\to0\), group the exact contact derivative before expanding:
\begin{align}
 h_\delta(\lambda^\star+s)
 ={}&\phi'_{\rm con}(\lambda^\star+s)
 +\frac{\delta^2}{2\beta s^2}(1+O(|s|)) \\[-0.2em]
 &+O(\delta^2)+O(\delta^4/s^4).
 \label{eq:boundary-balance-main}
\end{align}
Set \(s=-c\delta^{2/3}\) and divide by \(\delta^{2/3}\).  Every remainder
in \eqref{eq:boundary-balance-main} vanishes, and the limiting equation is
\[
 -\kappa c+\frac{1}{2\beta c^2}=0.
\]
Its positive root is \(c_0=(2\beta\kappa)^{-1/3}\).  Evaluating at fixed
\(c_-<c_0<c_+\) gives opposite signs for all small \(\delta\), so scalar
monotonicity brackets the unique root.  This also proves
\(|s_\delta|\asymp\delta^{2/3}\gg\delta\), validating the expansion.
Letting the brackets approach \(c_0\) proves \eqref{eq:boundary-rate-main}.
\end{proof}

The curvature condition is genuine.  For
\(A(\lambda^\star+s)=\diag(1+s,s)\), one has \(a=\beta=1\), but the contact
branch is constant on \(-1<s<0\), so \(\kappa=0\) and there is no isolated
two-thirds balance.

For the orientation \(a=\beta>0\) used in the crossover theorem below,
this contact curvature is exactly
\(\kappa=\ell''(0)-\tau''(0)=\kappa_0\).

For the boundary member of \cref{thm:open-2by2} with \(b_1=b_2=c=1\) and
\(d=0\), one has \(\lambda^\star=1\), \(a=\beta=1/2\), and
\(\phi'_{\rm con}(1+s)=s/\sqrt{s^2+4}=s/2+O(s^3)\) on the left.  Hence
\begin{equation}
 1-\lambda_\delta\sim2^{1/3}\delta^{2/3}.
 \label{eq:boundary-example-main}
\end{equation}

The exponent is not intrinsic to the certificate slice alone: it also records
how the chosen profile approaches the spectral-ball boundary.

\begin{corollary}[The boundary exponent is a tail law]
\label{cor:profile-tail-rate-main}
Under the hypotheses of \cref{thm:rate-boundary}, let
\(\lambda_\delta^\psi\) be any minimizer of the objective induced by an admissible
profile from \cref{prop:smoothing-family-main}.
\begin{enumerate}
 \item If, for some \(p>1\) and \(\gamma_\psi>0\),
 \begin{equation}
  \psi'(x)=1-\gamma_\psi x^{-p}+o(x^{-p})
  \qquad(x\to\infty),
  \label{eq:profile-polynomial-tail-main}
 \end{equation}
 then
 \begin{equation}
  \lambda_\delta^\psi-\lambda^\star
  =-\sign(a)
  \left(\frac{\gamma_\psi}{\kappa\beta^{p-1}}\right)^{1/(p+1)}
  \delta^{p/(p+1)}(1+o(1)).
  \label{eq:profile-tail-rate-main}
 \end{equation}
 \item If \(\psi'\) first reaches one at a finite threshold
 \(L_\psi>0\) and is identically one thereafter, then
 \begin{equation}
  \lambda_\delta^\psi-\lambda^\star
  =-\sign(a)\frac{L_\psi}{\beta}\,\delta+o(\delta).
  \label{eq:saturated-boundary-rate-main}
 \end{equation}
 For the standard Huber profile in \eqref{eq:two-profile-members-main},
 \(L_\psi=1\) and the sharper expansion is
 \begin{equation}
  \lambda_\delta^{\rm hub}-\lambda^\star
  =-\sign(a)\frac{\delta}{\beta}+O(\delta^2).
  \label{eq:huber-boundary-rate-main}
 \end{equation}
\end{enumerate}
Pseudo-Huber has \(p=2\) and \(\gamma_\psi=1/2\), so
\eqref{eq:profile-tail-rate-main} recovers
\eqref{eq:boundary-rate-main}; finite saturation yields the separate linear
boundary regime.
\end{corollary}

\begin{proof}
It is enough to take \(a=\beta>0\), for which the contact side is \(s<0\).
Along a scale with \(\delta/|s|\to0\), the separated positive cluster differs
from its polar limit by \(o(|s|)\) on the scale obtained below.  Using
\(\tau(s)=\beta s+O(s^2)\) and grouping against the exact contact derivative,
\eqref{eq:profile-polynomial-tail-main} gives
\[
 h_{\delta}^{\psi}(\lambda^\star+s)
 =\kappa s+
 \frac{\gamma_\psi\delta^p}{\beta^{p-1}|s|^p}
 +o\!\left(|s|+\frac{\delta^p}{|s|^p}\right).
\]
Putting \(s=-c\delta^{p/(p+1)}\) and bracketing the monotone root yields
\(\kappa c^{p+1}=\gamma_\psi/\beta^{p-1}\), proving the first claim.

If saturation first occurs at \(L_\psi\), the separated modes are already
saturated for all small \(\delta\).  Monotonicity brackets the root immediately
inside the transition edge \(|\tau(s)|=L_\psi\delta\), and therefore
\(|s|/(L_\psi\delta/\beta)\to1\), proving
\eqref{eq:saturated-boundary-rate-main}.  For standard Huber, inside
\(|\tau(s)|<\delta\) the small mode gives the sharper expansion
\[
 h_\delta^{\rm hub}(\lambda^\star+s)
 =\beta+\frac{\beta^2s}{\delta}
   +O(|s|+s^2/\delta).
\]
The monotone root therefore satisfies
\(s=-\delta/\beta+O(\delta^2)\).  Reflection handles \(a=-\beta\).
\end{proof}

\begin{corollary}[Selected-direction upper bounds]
\label{cor:direction-rates-main}
In fixed dimension and under the corresponding hypotheses of
\cref{thm:rate-regular,thm:direction-strict-general,thm:rate-boundary},
\begin{equation}
 \|\Phi_\delta-\Phi_R^\star\|_F
 =O(\delta^2),\qquad O(\delta),\qquad O(\delta^{2/3}),
 \label{eq:direction-rates-main}
\end{equation}
respectively.  The first two statements are upper bounds with the refinements
already stated.  The boundary rate is always exact under its hypotheses,
as the nonzero direction coefficient in
\cref{cor:direction-crossover-main} shows by taking a constant path.
The strict entry holds for arbitrary
rectangular shape and corank; \cref{thm:direction-strict-general} gives its
first-order constant and the sharper \(O(\delta^2)\) rate when \(a=0\).
\end{corollary}

\begin{proof}
On the separated cluster, the polar factor and smoothed spectral map are
analytic in the matrix argument and \(\delta^2\), so their variation is the
multiplier displacement plus \(O(\delta^2)\); this proves the regular entry.
The strict entry is \eqref{eq:strict-direction-first-main}.  At a boundary
kink, both the singular-vector movement and the coefficient
error are \(O(\delta^{2/3})\) because
\(|s|=O(\delta^{2/3})\) and
\(\delta^2/s^2=O(\delta^{2/3})\).
\end{proof}

\subsection{A local strict-to-boundary crossover}

The fixed-instance constants in
\cref{thm:direction-strict-general,prop:strict-constants-closed-main}
deteriorate as \(\beta-|a|\downarrow0\): \(\eta_R\) diverges and \(d_R\)
vanishes.  The result below is the uniform replacement in precisely that
degenerating direction, currently for a transverse square corank-one
crossing.  Thus the boundary law is not confined to one isolated parameter
value; a moving normal form controls the entire nearby transition.

\begin{lemma}[Moving crossing for general matrix normals]
\label{lem:moving-crossing-main}
Let \(A(\lambda,z)=G(z)+\lambda\Theta(z)\), \(z\in\mathbb R^p\), where
\(G\) and \(\Theta\) are real analytic and \(\Theta(z)\ne0\).  Suppose
the matrices are square, \(A(\lambda_0,0)\) has corank one, and
\begin{equation}
 u_0^\top\Theta(0)v_0\ne0.
 \label{eq:transverse-family-main}
\end{equation}
After shrinking the neighborhood, there are an analytic crossing
\(\lambda_{\rm sing}(z)\) and jointly analytic functions \(\tau,\ell,w\),
with \(s=\lambda-\lambda_{\rm sing}(z)\), such that
\begin{equation}
 \|A(\lambda_{\rm sing}(z)+s,z)\|_*
 =\ell(s,z)+|\tau(s,z)|,
 \qquad
 \tau(s,z)=s\sqrt{w(s,z)},\qquad w(s,z)>0.
 \label{eq:moving-normal-form-main}
\end{equation}
The sign can be chosen so that
\(\beta(z)=\partial_s\tau(0,z)>0\), and \(\ell\) is analytic even if
singular values inside the separated positive cluster cross.
\end{lemma}

\begin{proof}
At a corank-one matrix the adjugate is a nonzero multiple of
\(v_0u_0^\top\).  Condition \eqref{eq:transverse-family-main} makes
\(\partial_\lambda\det A(\lambda_0,0)\ne0\), so the analytic implicit
function theorem gives \(\lambda_{\rm sing}(z)\).  Set
\(\widetilde A(s,z)=A(\lambda_{\rm sing}(z)+s,z)\) and
\(B(s,z)=\widetilde A(s,z)^\top\widetilde A(s,z)\).  For each nearby \(z\),
\(B(0,z)\) has a simple zero eigenvalue separated from the rest; its analytic
continuation \(\mu(s,z)\) satisfies \(\mu(0,z)=0\).  Analytic division gives
\(\mu=s\mu_1\).  Since \(B\succeq0\), the scalar function
\(s\mapsto\mu(s,z)\) has a minimum at zero, so a second division gives
\(\mu=s^2w\).  Transversality makes \(w>0\) after shrinking.

For the remaining modes, a fixed contour around the positive cluster of the
symmetric dilation defines a jointly analytic Riesz projection.  The trace of
the dilation against this projection is the analytic eigenvalue sum
\(\ell(s,z)\), proving \eqref{eq:moving-normal-form-main}.
\end{proof}

\begin{theorem}[Unified cubic-root crossover]
\label{thm:crossover-main}
In \cref{lem:moving-crossing-main}, suppose at \(z=0\)
\begin{equation}
 a_0=\partial_s\ell(0,0)=\beta_0>0,
 \qquad
 \kappa_0=\partial_{ss}\ell(0,0)-\partial_{ss}\tau(0,0)>0.
 \label{eq:crossover-boundary-data-main}
\end{equation}
Let \(z=z_\delta\to0\) approach from the strict or boundary side and set
\begin{align*}
 a_\delta&=\partial_s\ell(0,z_\delta),
 &\beta_\delta&=\partial_s\tau(0,z_\delta),\\
 \Delta_\delta&=\beta_\delta-a_\delta\ge0,
 &\kappa_\delta&=\partial_{ss}\ell(0,z_\delta)
                    -\partial_{ss}\tau(0,z_\delta).
\end{align*}
Let \(\lambda_\delta^\star=\lambda_{\rm sing}(z_\delta)\), let
\(\lambda_\delta\) be the smoothed minimizer, and put
\(x_\delta=\lambda_\delta^\star-\lambda_\delta\).  Define
\(\widehat x_\delta>0\) as the unique positive root of
\begin{equation}
 \kappa_\delta\widehat x_\delta^3
 +\Delta_\delta\widehat x_\delta^2
 =\frac{\delta^2}{2\beta_\delta}.
 \label{eq:cubic-root-main}
\end{equation}
Then \(x_\delta>0\) for all small \(\delta\), and, along every such path,
including oscillatory ones,
\begin{equation}
 \frac{x_\delta}{\widehat x_\delta}
 =1+O(\Delta_\delta+\widehat x_\delta).
 \label{eq:cubic-relative-main}
\end{equation}
Consequently:
\begin{enumerate}
 \item if \(\Delta_\delta/\delta^{2/3}\to\zeta\in[0,\infty)\), then
 \begin{equation}
  \frac{x_\delta}{\delta^{2/3}}\to y_\zeta,
  \qquad
  \kappa_0y_\zeta^3+\zeta y_\zeta^2=\frac{1}{2\beta_0};
  \label{eq:crossover-critical-main}
 \end{equation}
 \item if \(\Delta_\delta/\delta^{2/3}\to\infty\), then
 \begin{equation}
  \frac{x_\delta\sqrt{\Delta_\delta}}{\delta}
  \to\frac{1}{\sqrt{2\beta_0}};
  \label{eq:crossover-inner-main}
 \end{equation}
 \item for arbitrary paths,
 \begin{equation}
  x_\delta\asymp\widehat x_\delta\asymp
  \min\left\{\frac{\delta}{\sqrt{\Delta_\delta}},
              \delta^{2/3}\right\},
  \label{eq:crossover-scale-main}
 \end{equation}
 where the first entry is \(+\infty\) at \(\Delta_\delta=0\).
\end{enumerate}
\end{theorem}

\begin{proof}
The one-sided derivatives at the moving crossing are
\(-\Delta_\delta\le0\) and
\(a_\delta+\beta_\delta\to2\beta_0>0\), so
\(\lambda_\delta^\star\) is an exact minimizer.  It is unique: if
\(\Delta_\delta>0\), the left derivative is strictly negative; if
\(\Delta_\delta=0\), positive contact curvature makes the derivative
strictly negative immediately to the left.  The derivative is strictly
positive immediately to the right in both cases, and convexity then
precludes any second minimizer.  For the instance indexed by
\(z\), set
\[
 h_{\delta,z}(\lambda)
 =\langle\Theta(z),Z_\delta(A(\lambda,z))\rangle,
 \quad
 \Delta(z)=\beta(z)-a(z),
 \quad
 \kappa(z)=\partial_{ss}\ell(0,z)-\partial_{ss}\tau(0,z),
\]
and \(H_{\delta,z}(x)=h_{\delta,z}(\lambda_{\rm sing}(z)-x)\).
Shrink a compact neighborhood of the reference point so that
\(\beta(z)\), \(\kappa(z)\), and every separated positive singular value
have fixed positive lower bounds there.  Joint analyticity gives bounded
derivatives on that neighborhood, and therefore the uniform expansions
\begin{align*}
 \ell_s(-x,z)-\tau_s(-x,z)&=-\Delta(z)-\kappa(z)x+O(x^2),\\
 \tau(-x,z)&=-\beta(z)x+O(x^2),\qquad
 \tau_s(-x,z)=\beta(z)+O(x).
\end{align*}
The separated smoothed cluster differs from \(\ell_s\) by \(O(\delta^2)\)
uniformly.  On the contact side, \(\tau(-x,z)<0\); its correction relative
to the unsmoothed small mode is exactly
\[
 \tau_s(-x,z)\left[1-
 \left(1+\frac{\delta^2}{\tau(-x,z)^2}\right)^{-1/2}\right]
 =\frac{\delta^2}{2\beta(z)x^2}
  +O\!\left(\frac{\delta^2}{x}+\frac{\delta^4}{x^4}\right).
\]
The binomial remainder is uniform when \(\delta/x\) is bounded by a
sufficiently small fixed constant.  Combining these bounds gives constants
\(z_0,x_0,\eta_0,C>0\), independent of the path, such that
whenever \(\|z\|\le z_0\), \(0<x\le x_0\), and
\(0<\delta\le\eta_0x\),
\begin{align}
 H_{\delta,z}(x)
 ={}&-\Delta(z)-\kappa(z)x
 +\frac{\delta^2}{2\beta(z)x^2}+E_{\delta,z}(x),
 \label{eq:crossover-functional-main}\\[-0.2em]
 |E_{\delta,z}(x)|
 \le{}&C\left(x^2+\frac{\delta^2}{x}
                    +\delta^2+\frac{\delta^4}{x^4}\right).
 \label{eq:crossover-remainder-main}
\end{align}
At \(x=0\), \(H_{\delta,z_\delta}(0)=a_\delta+O(\delta^2)>0\), whereas
it is negative at a sufficiently small fixed \(x_0>0\).  Strict convexity
places the unique root in \((0,x_0)\).

Put
\(M_\delta=\Delta_\delta+\kappa_\delta\widehat x_\delta\).  Equation
\eqref{eq:cubic-root-main} is equivalent to
\[
 \frac{\delta^2}{2\beta_\delta\widehat x_\delta^2}=M_\delta.
\]
It implies \(\widehat x_\delta,M_\delta\to0\) and
\(\delta/\widehat x_\delta=\sqrt{2\beta_\delta M_\delta}\to0\), so the
uniform expansion applies at \(x=c\widehat x_\delta\).  Uniformly for \(c\)
in a compact subset of \((0,\infty)\), the four terms in the remainder,
divided by \(M_\delta\), are respectively
\[
 O(\widehat x_\delta),\quad O(\widehat x_\delta),\quad
 O(\widehat x_\delta^2),\quad O(M_\delta).
\]
The leading part at this point is
\[
 \Delta_\delta(c^{-2}-1)
 +\kappa_\delta\widehat x_\delta(c^{-2}-c).
\]
It is a fixed positive multiple of \(M_\delta\) for fixed \(c<1\), and a
fixed negative multiple for fixed \(c>1\).  Monotonicity therefore gives
\(x_\delta/\widehat x_\delta\to1\).  More quantitatively, for
\(c=1+u\) with \(|u|\) small, the leading term has the sign of \(-u\) and
magnitude at least \(c_0|u|M_\delta\), whereas the remainder is bounded by
\(C_1(\Delta_\delta+\widehat x_\delta)M_\delta\), with uniform positive
constants \(c_0,C_1\).  Taking
\(u=\pm K(\Delta_\delta+\widehat x_\delta)\) and
\(K>C_1/c_0\) gives the quantitative bracket
\eqref{eq:cubic-relative-main}.

The three consequences now follow directly from
\eqref{eq:cubic-root-main}: divide it by \(\delta^2\) in the critical window,
drop its cubic term when \(\Delta_\delta/\delta^{2/3}\to\infty\), and compare
the two nonnegative terms for the two-sided order bound.
\end{proof}

The cubic also controls the matrix actually returned by the oracle.  The
comparator must be the center of the moving instance rather than the fixed
center at \(z=0\).

\begin{corollary}[Direction-level crossover]
\label{cor:direction-crossover-main}
Under \cref{thm:crossover-main}, let \(P(s,z)\) be the analytic polar
contribution of the separated positive singular cluster and let \(D(s,z)\) be
the analytic signed dyad of the small mode, chosen so that
\(\langle\Theta(z),D(0,z)\rangle=\beta(z)>0\).  With derivatives evaluated at
\((s,z)=(0,0)\), define
 \begin{equation}
 \mathcal M_0=-\partial_sP(0,0)+\partial_sD(0,0)
 +\frac{\kappa_0}{\beta_0}D(0,0).
 \label{eq:direction-M0-main}
\end{equation}
Then, relative to the exact Fenchel center \(\Phi_R^\star(z_\delta)\) of the
moving instance,
\begin{equation}
 \frac{\Phi_\delta-\Phi_R^\star(z_\delta)}{x_\delta}
 \longrightarrow\mathcal M_0.
 \label{eq:direction-crossover-limit-main}
\end{equation}
Moreover,
\begin{equation}
 \langle D(0,0),\mathcal M_0\rangle=\frac{\kappa_0}{\beta_0}>0,
 \qquad
 \|\mathcal M_0\|_F\ge\frac{\kappa_0}{\beta_0}.
 \label{eq:direction-M0-nonzero-main}
\end{equation}
Consequently, along every path in the theorem,
\[
 \|\Phi_\delta-\Phi_R^\star(z_\delta)\|_F
 \sim\|\mathcal M_0\|_F\widehat x_\delta.
\]
In particular, the direction-level crossover cannot lose its leading term
under the printed positive-curvature assumption.
\end{corollary}

\begin{proof}
The smoothed spectral decomposition at \(s=-x_\delta\) is
\[
 \Phi_\delta=P_\delta(-x_\delta,z_\delta)
 +r_\delta(-x_\delta,z_\delta)D(-x_\delta,z_\delta),
 \quad
 r_\delta=\frac{\tau}{\sqrt{\tau^2+\delta^2}},
\]
where \(P_\delta=P+O(\delta^2)\).  Tangency gives
\(r_\delta=-\ell_\delta'/\tau'\).  Expanding this quotient and using
\(a=\beta-\Delta\) and
\(\kappa=\ell_{ss}-\tau_{ss}\) gives, uniformly,
\[
 r_\delta(-x,z)+\frac{a(z)}{\beta(z)}
 =x\left\{\frac{\kappa(z)}{\beta(z)}
 +\frac{\Delta(z)\tau_{ss}(0,z)}{\beta(z)^2}\right\}
 +O(x^2+\delta^2).
\]
Since \(\Delta(z_\delta)\to0\) and \(\delta^2/x_\delta\to0\), this yields
\[
 \frac{r_\delta(-x_\delta,z_\delta)+a(z_\delta)/\beta(z_\delta)}
 {x_\delta}\longrightarrow\frac{\kappa_0}{\beta_0}.
\]
Since
\(\Phi_R^\star(z)=P(0,z)-[a(z)/\beta(z)]D(0,z)\), Taylor expansion proves
\eqref{eq:direction-crossover-limit-main}; the separated-cluster smoothing
error is negligible because \(\delta^2/x_\delta\to0\).
To show that the limit is nonzero, write \(D=u_0v_0^\top\) in compatible
analytic unit frames.  Its Frobenius norm is one, so
\(\langle D,D_s\rangle=0\).  Also
\(P=U_+Q V_+^\top\), where \(Q\) is the positive-cluster polar factor.
Differentiation and \(U_+^\top u_0=V_+^\top v_0=0\) give
\(\langle D,P_s\rangle=0\): each of the three differentiated terms
retains at least one of these zero factors.  Taking the pairing of
\eqref{eq:direction-M0-main} with \(D(0,0)\) proves
\eqref{eq:direction-M0-nonzero-main}, including its norm bound by
Cauchy--Schwarz.  The asymptotic norm formula follows from
\eqref{eq:cubic-relative-main}.
\end{proof}

The moving-crossing recentering is essential: nonnegativity together with a
single identity \(\mu(0,0)=0\) would not imply the factorization
\(\mu=s^2w\).  The theorem is uniform only near one fixed analytic transverse
corank-one boundary instance.
